\documentclass[11pt,a4paper,fontset=fandol,scheme=plain]{ctexart}
\usepackage[zh]{math-writing}
\begin{document}
\mwtitle{文章标题}{作者 · 日期 · 版本}
\section{主题标题}
% 先在正文写明已有背景、所用约定和本节的目标。
\begin{definition}[概念名称]\label{def:main}
在这里写完整定义，说明对象、数据与条件。
\end{definition}
\begin{theorem}[结果名称]\label{thm:main}
在这里写假设与结论；不要用“见上文”掩盖实质条件。
\end{theorem}
\begin{proof}
在这里写证明。长证明可先给化归，短证明直接进入论证。
\end{proof}
\begin{remark}\label{rem:scope}
说明结论的用途或边界；不需要时删除本段。
\end{remark}
引用示例：\cref{def:main,thm:main,rem:scope}。
\end{document}
